% ch13 — Initial Data
\chapter{Initial Data}
\label{ch:initial-data}

A valid initial-data set must satisfy the Hamiltonian and momentum constraints
on the initial hypersurface. This chapter develops the conformal decomposition
that converts the constraint equations into solvable elliptic PDEs and constructs
the black-hole initial data used in the simulations.

\section{The Constraint Equations}
\label{sec:id-constraints}

We state the Hamiltonian and momentum constraints as elliptic equations on the
initial slice and explain why free evolution schemes require they be solved
exactly at $t=0$.

\section{Conformal Flatness and the Lichnerowicz Equation}
\label{sec:id-conformal-flat}
\coderef{initial\_data}{rs}

We adopt a conformally flat spatial metric and derive the Lichnerowicz equation
for the conformal factor $\psi$, reducing the Hamiltonian constraint to a
single nonlinear elliptic PDE.

\section{Isotropic Schwarzschild and Brill-Lindquist Data}
\label{sec:id-brill-lindquist}

We construct isotropic Schwarzschild initial data for a single black hole,
then generalise to the Brill-Lindquist superposition for multiple momentarily
stationary holes.

\section{Bowen-York Data}
\label{sec:id-bowen-york}

We add linear and angular momentum via the Bowen-York extrinsic curvature
ansatz and solve the resulting coupled system for spinning and boosted
black holes.

\section{XCTS Formulation\starmark{}}
\label{sec:id-xcts}
\coderef{xcts}{rs}

We present the extended conformal thin-sandwich decomposition, which provides
greater physical control over the initial slice, and outline the GMRES-based
solver used in the codebase.
\coderef{gmres}{rs}

\paragraph{Key References.}
The conformal decomposition follows \cite{York1979,Cook2000}; Bowen-York data
follows \cite{Bowen1980}; XCTS follows \cite{Pfeiffer2003}; the solver
strategy draws on \cite{Baumgarte2010}.
